\documentclass[11pt]{article}
\usepackage{microtype,amsmath,amssymb,amsthm,mathtools,booktabs,array}
\usepackage[margin=1.05in]{geometry}
\usepackage[colorlinks=true,linkcolor=blue!45!black,citecolor=blue!45!black,urlcolor=blue!45!black]{hyperref}
\usepackage{xcolor}
\hypersetup{pdftitle={Constant quadratic Gaussian fibres and repeated-line cubics},pdfsubject={Working mathematical manuscript}}
\newtheorem{theorem}{Theorem}[section]
\newtheorem{proposition}[theorem]{Proposition}
\newtheorem{lemma}[theorem]{Lemma}
\newtheorem{corollary}[theorem]{Corollary}
\theoremstyle{definition}\newtheorem{definition}[theorem]{Definition}
\theoremstyle{remark}\newtheorem{remark}[theorem]{Remark}
\newcommand{\E}{\mathbb E}
\newcommand{\C}{\mathbb C}
\newcommand{\tr}{\operatorname{tr}}
\newcommand{\wt}{\operatorname{wt}}
\newcommand{\G}{\mathsf G}
\DeclareMathOperator{\Span}{span}
\title{Constant quadratic Gaussian fibres and repeated-line cubics}
\author{}\date{17 November 2026: internal written proof}
\begin{document}\maketitle
This note proves an all-degree Gaussian fibre theorem and the repeated-line
cubic reduction. It includes the full valued-field polynomial stationary
criterion used in its proof. All original markers and determinant
normalizations are retained. It uses no full image-kernel computation.
Classical formal stationary mechanisms are credited; arbitrary formal
Laurent coefficients are allowed in the polynomial frame theorem, but
no arbitrary-ring descent is inferred. No novelty, independent replay,
formal verification or external review is claimed.

For independent standard Gaussian coordinates extend expectation
complex-linearly and define
\[ \Phi_Q(t)=\sum_{j\ge0}\mathbb E[QP^j]t^j/j!. \]
Call $P$ pure if $\Phi_1=1$, and safe if $\Phi_Q$ is polynomial
for every original polynomial marker $Q$. In ternary circular
coordinates the Gaussian phase is $ZW+T^2/2$, with
$\mathbb E[Z^aW^b]=\delta_{ab}a!$.
The top quotient has spatial relations $\partial_i+tP_i-x_i$ and
parameter connection $\partial_t+P$. Coefficientwise integration by
parts defines its formal Gaussian functional.
\begin{lemma}[Rational seed restriction]\label{lem:seedfunctional}
A horizontal formal functional taking a seed to one restricts to a
rational horizontal quotient on its differential seed submodule.
\end{lemma}
\begin{proof}
Its values on the seed and all successive derivatives are $1,0,0,\ldots$.
Every rational relation is respected, so the restriction is a well-defined
map to $\mathbb C(t)$. Exact field extension preserves it.
\end{proof}
\begin{lemma}[The algebraic frame test]\label{lem:frame}
There are Gaussian coordinates $(Z,W,T)$ in which $P$ is affine linear
in $W$ if and only if there is a nonzero vector $v\in\C^3$ such that
\begin{equation}\label{eq:frame}
 v^Tv=0,\qquad \nabla F_3(v)=0,\qquad F_2(v)=0.
\end{equation}
\end{lemma}
\begin{proof}
Polarization gives
\[
 D_v^2P(x)=2\nabla F_3(v)^Tx+2F_2(v).
\]
The last two conditions in \eqref{eq:frame} say exactly that this
polynomial vanishes. Complete the nonzero null vector $v$ to a Witt
basis: choose a null partner with pairing one, and a unit vector in
their orthogonal complement. In these coordinates $D_v$ is the $W$
derivative. Hence $D_v^2P=0$ says that no second or higher power of
$W$ occurs. The converse reads off $v$ from the $W$ direction.
\end{proof}

\subsection{A formal leading-coefficient criterion}\label{sec:stationaryballs}

For a polynomial phase, the remaining shrinking argument can be
replaced by a valuation calculation. The construction below includes
critical points that collide, remain multiple, or escape to infinity.
The key is to group them before making a Gaussian approximation.

Write $v$ for the valuation of a formal parameter $r$, normalized by
$v(r)=1$. The valuation of a polynomial is the least valuation of its
coefficients. We use the formal Newton--Puiseux theorem: finitely many
algebraic elements over $\C((r))$ belong to one field
$\C((r^{1/N}))$. Convergence is not required~\cite[Chapter~2]{Wall}.

\begin{lemma}[Polynomial separation of formal balls]\label{lem:ballseparation}
Let finitely many disjoint balls have centers $c_B$ in a Puiseux field
and radii specified by $v(x-c_B)\geq\rho_B$. After one common
ramification, put $\delta_B=r^{\rho_B}$. For any one ball $B$ and
integer $j\geq0$, there is a polynomial $q_{B,j}(x)$ over the Laurent
field whose substitution $x=c_C+\delta_C y$ has integral coefficients
for every ball $C$, and whose reduction is $y^j$ on $B$ and zero on
every other ball.
\end{lemma}
\begin{proof}
The finitely many valuations $v(c_B-c_C)$ give a finite tree of
clusters of centers. At a node with scale $r^d$, the coordinate
$u=(x-c)/r^d$ has distinct constant reductions $a_1,\ldots,a_s$
on its children. The Lagrange polynomial
\[
 e_i(u)=\prod_{k\ne i}\frac{u-a_k}{a_i-a_k}
\]
reduces to one on child $i$ and zero on every other child.

Start at the chosen leaf with $((x-c_B)/\delta_B)^j$. Induct upward.
Suppose a polynomial $h$ already has the required reductions on the
chosen child. On each of the other child balls its coefficients have
some finite lower valuation. The restriction of $e_i$ there has
strictly positive valuation. Consequently $e_i^M h$, for a sufficiently
large fixed integer $M$, has positive coefficient valuation on all
those children. On the chosen child $e_i^M=1+$ positive-order terms,
so the existing reductions are unchanged. There are finitely many
nodes, and this constructs the required polynomial. A node with just
one leaf requires no filter. This also explains why independently
normalized Lagrange polynomials, with no control of their growth on
other scales, would not suffice.
\end{proof}

\begin{theorem}[Formal stationary balls]\label{thm:stationaryballs}
Let $\varphi(r,x)\in\C((r))[x]$ have degree $D\geq2$, and let
$h(r)\in\C((r))\setminus\{0\}$. Put
\[
 M=\C((r))[x]/(\partial_x+\varphi_x)\C((r))[x],\qquad
 \nabla[q]=[q_r+(\varphi_r+h'/h)q].
\]
After a finite ramification, this module has a full frame of rank-one
formal exponential functionals. If $a_D$ is the leading coefficient
of $\varphi$, every zero-exponential functional in this frame takes
the seed $[1]$ to a series of valuation at least
\begin{equation}\label{eq:valuationbound}
 v(h)-\frac{v(a_D)}D.
\end{equation}
In particular, if $Dv(h)>v(a_D)$, the differential submodule generated
by $[1]$ has no horizontal quotient taking $[1]$ to the constant one.
\end{theorem}
\begin{proof}
The two twisted derivatives commute, so the parameter connection
is well defined on the quotient. A nonzero polynomial of degree $m$
has spatial twisted derivative of degree $m+D-1$. Reduction of
successive leading terms therefore gives a basis
$1,x,\ldots,x^{D-2}$ and dimension $D-1$.

Split $\varphi_x$, retaining root multiplicities. For a critical root
$\xi$, expand
\[
 \varphi(\xi+u)=c_\xi+\sum_{j=2}^D a_{\xi,j}u^j,
 \qquad
 \rho_\xi=\max_{2\leq j\leq D,\ a_{\xi,j}\ne0}
                    \frac{-v(a_{\xi,j})}{j}.
\]
Enlarge the ramification so that all these scales are integral, and
set $\delta_\xi=r^{\rho_\xi}$. The polynomial
\[
 \varphi(\xi+\delta_\xi y)-c_\xi
\]
has integral coefficients, no linear term, and at least one unit
coefficient. Its reduction $F_\xi(y)$ has degree $d_\xi\geq2$.
Here and below we may replace $r$ by a root of $r$ and renormalize
the valuation; every displayed strict inequality is unchanged.
As $a_{\xi,D}=a_D$,
\begin{equation}\label{eq:balllower}
 \rho_\xi\geq-v(a_D)/D.
\end{equation}

\emph{The balls partition all critical roots.}
If another critical root $\zeta$ satisfies
$v(\zeta-\xi)\geq\rho_\xi$, translate the integral scaled polynomial
by $(\zeta-\xi)/\delta_\xi$. All its coefficients remain integral,
and the coefficient of degree $d_\xi$ remains a unit: higher-degree
terms have positive valuation. The linear term vanishes because
$\zeta$ is critical. It follows from the definition of the maximum
that $\rho_\zeta=\rho_\xi$. Thus their balls are identical.
If two of the balls intersect, one contains the other; the same
argument makes them identical. Otherwise they are disjoint.

The scaled derivative has reduction $F_\xi'$ of degree $d_\xi-1$.
It therefore has exactly $d_\xi-1$ integral roots, with multiplicity.
Here is a direct check of that count: factor it in a Puiseux field.
An integral root contributes a linear factor to the reduction;
for a root of negative valuation, factor out that root and the
remaining factor reduces to a constant. Multiplicativity of the
coefficient valuation leaves a nonzero constant scalar. Hence the
degree of the reduction counts exactly the integral roots.
The ball $B$ thus contains $d_B-1$ critical roots. Since
$\deg\varphi_x=D-1$,
\begin{equation}\label{eq:ballrankcount}
 \sum_B(d_B-1)=D-1.
\end{equation}

\emph{A functional for every dimension in a ball.}
The fixed polynomial $F_B$ has twisted quotient
\[
 H_B=\C[y]/(\partial_y+F_B')\C[y]
\]
with basis $1,y,\ldots,y^{d_B-2}$. Let $\ell^0_{B,j}$ be its dual
basis. Write
\[
 \varphi(\xi_B+\delta_B y)=c_B+F_B(y)+E_B(r,y),
 \qquad E_B\in r\C[[r]][y].
\]
For a polynomial amplitude $q$ define
\begin{equation}\label{eq:ballfunctional}
 \ell_{B,j}(q)=h\delta_B e^{c_B}
   \ell^0_{B,j}\!\left(e^{E_B(r,y)}q(\xi_B+\delta_B y)\right).
\end{equation}
Each coefficient uses finitely many polynomial terms. The exponential
in this formula is purely formal. The product rule conjugates the
spatial twisted derivative to $\partial_y+F_B'$, so the functional
annihilates the defining relations of $M$.

We spell out the parameter check because it retains the moving volume.
At fixed $y$, put $w=\xi_B'+\delta_B'y$. Differentiating
\eqref{eq:ballfunctional}, then subtracting its value on
$q_r+(\varphi_r+h'/h)q$, leaves, inside the fixed functional,
\[
 e^{E_B}\left(wq_x+\frac{\delta_B'}{\delta_B}q+\varphi_xwq\right)
 =\frac1{\delta_B}(\partial_y+F_B')
                      \left(e^{E_B}wq\right).
\]
It is zero. Thus the functionals are horizontal, including the
Jacobian $\delta_B$ and normalization $h$.

\emph{Completeness, not just a rank count.}
Choose the polynomials $q_{B,j}$ of
Lemma~\ref{lem:ballseparation}, for $0\leq j<d_B-1$.
Divide each evaluation column by its scalar $h\delta_Be^{c_B}$.
The resulting matrix reduces to the identity: $E_B$ has positive
order and the reductions of the amplitudes are the chosen monomials
on one ball and zero elsewhere. The matrix is therefore invertible.
Leading-degree reduction gives $\dim M=D-1$, so
\eqref{eq:ballrankcount} accounts for the entire module.

Remove from $c_B$ its finite negative-power part $q_B(r)$.
The remaining exponential is an invertible formal power series,
up to a nonzero constant. Formula~\eqref{eq:ballfunctional} and the
invertible matrix give the asserted direct sum of rank-one factors
$e^{q_B}$ times powers of $r$. If $q_B=0$, its value on the seed
has valuation at least $v(h)+\rho_B$, proving
\eqref{eq:valuationbound}.

Finally, after ramification this is a semisimple differential module.
A quotient on the submodule generated by the seed extends to the
whole module. Nonzero exponential factors have no horizontal map
to the Laurent field, because the logarithmic derivative of a
Laurent series cannot be the derivative of a nonzero negative-power
exponent. On the zero part, every horizontal functional is a
constant combination of those in the full frame. If the bound is
strictly positive, their seed values cannot sum to one.
\end{proof}

The construction is a formal stationary-phase argument. Its classical
background includes Newton--Puiseux and the formal stationary-phase
theory of exponential connections; Sabbah's local Fourier--Laplace
formula is a primary reference for the latter~\cite{SabbahStationary}.
We do not identify that formula's scope with the particular proof
above or claim priority for this formulation. The proof here also
allows nonconvergent formal Laurent coefficients. It uses a valued
field, not an arbitrary coefficient ring.

\begin{corollary}[Leading-coefficient obstruction]\label{cor:leadingobstruction}
For a rational polynomial phase $\varphi(t,x)$ of degree $D\geq2$
and nonzero rational scalar normalization $h(t)$, let $v_\infty$
mean order at $\varepsilon=1/t=0$. If
\begin{equation}\label{eq:leadingcriterion}
 v_\infty(h^D/a_D)>0,
\end{equation}
its seed submodule cannot have a rational horizontal quotient with
seed value one. In particular any horizontal formal functional at
another point cannot take that seed to one.
\end{corollary}
\begin{proof}
Apply the theorem at infinity. A rational horizontal quotient remains
a quotient on extension to the formal Laurent field. A formal
functional with seed value one restricts to such a rational quotient:
its values on successive seed derivatives are $1,0,0,\ldots$.
\end{proof}

The inequality must be strict. For example
$\varphi=-x^2/(2t^2)$ and $h=1/t$ have
$v_\infty(h^2/a_2)=0$, and their normalized Gaussian period is a
nonzero constant. Thus the earlier adverse scaling example is also
an exact boundary test of the criterion.

\subsection{Quadratic elimination and the repeated lines}
\label{sec:repeatedlines}

The final two leading types have a repeated factor. They reduce
to a polynomial in one variable after eliminating two Gaussian
variables. The full leading-coefficient theorem makes an escape
table unnecessary: only the actual degree, its leading
coefficient and the determinant factor are needed.

\begin{lemma}[The determinant factor in quadratic elimination]
\label{lem:quadraticelimination}
Over $\C(t)$, let a phase have the form
\[
 -\tfrac12 y^T M(t)y+L(t,x)^Ty+V(t,x),
\]
where $M$ is an invertible symmetric matrix and $L,V$ are
polynomial in $x$. Eliminating $y$ identifies its top twisted
module and constant volume seed with the one-variable phase
and normalization
\[
 \varphi=V+\tfrac12L^TM^{-1}L,\qquad
 h=(\det M)^{-1/2}.
\]
The parameter connection is
$\partial_t+\varphi_t+h'/h$. An algebraic square root is allowed.
If the actual degree $d$ of $\varphi$ is at least two and
$d\,v_\infty(h)>v_\infty([x^d]\varphi)$, its constant seed
cannot have a rational unit quotient.
\end{lemma}
\begin{proof}
The change $y=z+M^{-1}L$ has Jacobian one. Its spatial twisted
relations reduce the $z$ factor to the one-dimensional Gaussian
quotient, with normalized contractions
$\ell(z_i z_j)=(M^{-1})_{ij}$. Parameter differentiation of its
constant class inserts $-z^TM'z/2$, hence
\[
 -\tfrac12\operatorname{tr}(M^{-1}M')=h'/h.
\]
Terms from the moving center are spatial exact forms, by the
product rule or Cartan's formula. This proves the identification
of both connection and seed. After a finite cover at infinity,
$h$ is a Laurent series and Theorem~\ref{thm:stationaryballs}
applies. The cover multiplies both sides of the valuation
inequality by the same positive integer.

The elimination is over the generic field. It need not define
a conditional formal average at $t=0$. The original Wick
functional first gives its rational seed quotient by
Lemma~\ref{lem:seedfunctional}; this quotient is transported
through the displayed module identification. No image-kernel
calculation or convergent original period is needed.
\end{proof}

\begin{lemma}\label{lem:quadraticfibrehighdegree}
Let $u,Y_1,\ldots,Y_s$ be independent standard Gaussian variables
and
\[
 P=Y^TAY+B(u)^TY+g(u),\qquad A=A^T,
\]
with $B,g$ arbitrary polynomials. Put
\[
 D(t)=\det(I-2tA),\quad C(t)=(I-2tA)^{-1},\quad
 \varphi=tg-u^2/2+(t^2/2)B^TCB.
\]
If $\deg_u\varphi\ge3$, then $P$ is not pure.
\end{lemma}
\begin{proof}
Quadratic elimination gives $h=D^{-1/2}$. Let $d=\deg D$ and
$r=\deg_u\varphi\ge3$. Its leading coefficient is $N(t)/D(t)$,
where $N$ is a nonzero polynomial divisible by $t$:
the degree-$r$ coefficient vanishes at $t=0$. Consequently
\[
 r v_\infty(h)-v_\infty(N/D)
     =(r/2-1)d+\deg N>0.
\]
Apply Lemma~\ref{lem:quadraticelimination}. This estimate is
independent of the degrees of $B,g$ and the dimension $s$.
\end{proof}

\subsubsection*{A nonisotropic repeated factor}

Normalize the repeated factor to $u$ with $\E[u^2]=1$ and take
an orthogonal Gaussian complement $Y\in\C^2$. Every such cubic
has the form in Lemma~\ref{lem:quadraticfibrehighdegree}, with
\[
 B(u)=p u^2+q u+r,\qquad\deg g\le3.
\]
For a double line $p\ne0$. If its effective quartic coefficient
is nonzero, the lemma excludes purity immediately. If that
coefficient vanishes identically, then
\[
 0=p^T\operatorname{adj}(I-2tA)p
   =p^Tp-2t\bigl(\operatorname{tr}(A)p^Tp-p^TAp\bigr).
\]
Thus $p^Tp=p^TAp=0$. The nonzero vector $(0,p)$ is isotropic,
lies in the leading singular plane $u=0$, and satisfies
$F_2(0,p)=0$. It supplies the frame of Lemma~\ref{lem:frame}.
For a triple line, $p=0$ and $[u^3]g\ne0$. The effective
phase has exact cubic coefficient $t[u^3]g$, so the same
high-degree lemma excludes purity. This exhausts all matrices
$A$, including nilpotent and defective ones, without a matrix
case split.

\subsubsection*{An isotropic repeated factor}

Use Gaussian circular coordinates with repeated factor $Z$.
The singular plane is $Z=0$, whose only isotropic direction is
the $W$ axis. Write $\lambda=[W^2]F_2$. If $\lambda=0$, the
frame test applies. Assume $\lambda\ne0$.

For a double line write its other factor as
$aW+bT+cZ$. The Gaussian null rotation
\[
 Z=Z',\quad T=T'+sZ',\quad W=W'-sT'-s^2Z'/2
\]
removes the $T$ term if $a\ne0$, or removes the $Z$ term if
$a=0,b\ne0$. Scaling $P$ gives exactly these leading forms:
\[
 \mathrm I: Z^2W+a_0Z^3,\qquad
 \mathrm {II}:Z^2T,\qquad \mathrm {III}:Z^3.
\]
The last is the triple line. The rotation preserves the Gaussian
and keeps the $W^2$ coefficient nonzero. All lower terms remain
arbitrary.

Write $P=\lambda W^2+W H(Z,T)+G(Z,T)$. In type I,
$H=Z^2+bZ+cT+d$ and $G=a_0Z^3+G_2$. In types II and III,
$H=bZ+cT+d$ and $G=Z^2T+G_2$ or $Z^3+G_2$ respectively.
Here $G_2$ is an arbitrary polynomial of degree at most two.
Eliminating $W$ in $tP-ZW-T^2/2$ gives
\[
 \Psi=t\left(G-\frac{H^2}{4\lambda}\right)
       +\frac{ZH}{2\lambda}-\frac{T^2}{2}
       -\frac{Z^2}{4\lambda t}.
\]
Its $T^2$ coefficient is $C_T=a t-1/2$, where
$a=g_{02}-c^2/(4\lambda)$ and $g_{02}=[T^2]G_2$.
Eliminate $T$ as well. If $e=\deg_t C_T\in\{0,1\}$, the total
normalization has $v_\infty(h)=(1+e)/2$. The actual remaining
leading coefficients are
\[
 \begin{array}{c|c|c}
 \text{type}&[Z^d]\varphi&d\\\hline
 \mathrm I&\displaystyle
  -\frac{t}{4\lambda}-\frac{c^2t^2}{16\lambda^2C_T}
   =-\frac{t(2g_{02}t-1)}{8\lambda C_T}&4\\[5pt]
 \mathrm {II}&\displaystyle-\frac{t^2}{4C_T}&4\\[4pt]
 \mathrm {III}&t&3.
 \end{array}
\]
None is identically zero. Their valuations are at most $e-1$,
$e-2$, and $-1$, respectively. Each is strictly less than
$d(1+e)/2$. Lemma~\ref{lem:quadraticelimination} excludes
purity in all three types. In particular this includes the
boundary where the quartic coefficient tends to a nonzero
constant while lower coefficients grow.

\begin{proposition}\label{prop:repeatedcomplete}
Every pure ternary cubic with repeated leading factor has a
Gaussian affine-linear coordinate.
\end{proposition}
\begin{proof}
The nonisotropic and isotropic alternatives above exhaust the
repeated factor. In every surviving branch the explicit frame
test succeeds. Apply the algebraic frame lemma.
\end{proof}

\section{An all-degree consequence of quadratic elimination}
\label{sec:constantfibres}

The estimate used for the repeated lines gives a useful result
in arbitrary dimension. It also identifies a precise remaining
problem about rational marked periods. This section is not
needed for the dimension--degree boundary.

\begin{theorem}\label{thm:constantfibres}
Let $u,Y_1,\ldots,Y_s$ be independent standard Gaussian variables,
$s\ge1$, and
\[
 P=Y^TAY+B(u)^TY+g(u),\qquad A=A^T,
\]
where $B,g$ have arbitrary degrees. Put $D=\det(I-2tA)$ and
$C=(I-2tA)^{-1}$. Purity is equivalent to the following finite
rational identities:
\begin{align}
 \deg g\le2,\qquad &\deg_u(B^TCB)\le2,\label{eq:fibredegree}\\
 D K=1,\qquad &\gamma+\beta^2/(2K)=0,\label{eq:fibreunit}
\end{align}
where $\alpha u^2+\beta u+\gamma=tg+(t^2/2)B^TCB$ and
$K=1-2\alpha$.

For every pure member, every original marked period is rational,
with possible finite poles only at zeros of $D$. If every
nonzero eigenvalue of $A$ is semisimple, all those poles cancel:
$P$ is safe. More explicitly, with
$b=\max(1,\deg_u B)$, interpreting $b=1$ if $B=0$,
\begin{equation}\label{eq:fibrecutoff}
 \deg_t\Phi_Q\le\bigl(s+(s+1)b\bigr)\deg Q.
\end{equation}
No restriction is imposed on the nilpotent zero-eigenspace.
The bound is sufficient, not asserted optimal.
\end{theorem}
\begin{proof}
Lemma~\ref{lem:quadraticfibrehighdegree} forces the effective
phase to be quadratic. Its higher-degree coefficients say
$tg_j+(t^2/2)[u^j]B^TCB=0$. At $t=0$ these first give $g_j=0$
for $j\ge3$, then the second identity in
\eqref{eq:fibredegree}. The unit series is exactly
\[
 \Phi_1=(DK)^{-1/2}\exp\!\left(\gamma+\frac{\beta^2}{2K}\right).
\]
Squaring the purity identity gives
$DK=\exp(2\gamma+\beta^2/K)$. A rational derivative has zero
residues, whereas the logarithmic derivative of a rational
function has its zero and pole orders as residues. Thus $DK$
is constant, and then so is $2\gamma+\beta^2/K$.
Normalization at zero gives
\eqref{eq:fibreunit}. Conversely those identities give
$\Phi_1=1$.

Set $\mu=D\beta$. It is polynomial because the only denominator
in $C$ is $D$. Under the normalized tilted functional the base
variable $u$ has mean $\mu$ and variance $D$. Conditional on
$u$, the $Y$ mean is $tCB(u)$ and its covariance is $C$.
Wick's rule therefore expresses every original marked period
as a rational function whose denominator is a power of $D$.
The unit period is one, so no additional exponential factor
remains. In particular their span over $\C(t)$ has rank one;
this assertion alone does not prove safety.

Assume now that all nonzero eigenvalues are semisimple. The
generalized eigenspaces of the symmetric matrix are orthogonal
and nondegenerate. Write $A_0$ for its nilpotent zero block.
Its resolvent $C_0$ is polynomial. Partial fractions in $t$
separate its contribution from the simple-pole terms at
nonzero eigenvalues. Thus each corresponding quadratic norm
in $B(u)$ has degree at most two in $u$.

Equation $K=1/D$ forces every nonzero eigenvalue to have
multiplicity one: the left side has at most a simple pole
there, while $1/D$ has its full eigenvalue multiplicity.
For such an eigenvalue $\lambda$, put
\[
 l_\lambda=1-2\lambda t,\quad t_\lambda=(2\lambda)^{-1},
 \quad D_\lambda=D/l_\lambda.
\]
Its scalar source is affine, $B_\lambda(u)=b_\lambda u+c_\lambda$,
because its square has degree at most two. Evaluating $DK=1$
at $t_\lambda$ gives the exact certificate
\begin{equation}\label{eq:fibrepolecertificate}
 t_\lambda^2b_\lambda^2D_\lambda(t_\lambda)=-1,
 \qquad
 \mu(t_\lambda)=-c_\lambda/b_\lambda.
\end{equation}
The second identity follows by evaluating $D\beta$ there.

The joint tilted functional of $u$ and all these nonzero-eigenvalue
coordinates is Gaussian. Its means and covariances are
\begin{align*}
 \E_tu&=\mu,&\E_tY_\lambda
   &=t(b_\lambda\mu+c_\lambda)/l_\lambda,\\
 \operatorname{Var}_t(u)&=D,&
 \operatorname{Cov}_t(u,Y_\lambda)&=tDb_\lambda/l_\lambda,\\
 \operatorname{Cov}_t(Y_\lambda,Y_\nu)
   &=t^2b_\lambda b_\nu D/(l_\lambda l_\nu)\quad(\lambda\ne\nu),&&\\
 \operatorname{Var}_t(Y_\lambda)
   &=(1+t^2b_\lambda^2D_\lambda)/l_\lambda.&&
\end{align*}
All are polynomials. The off-diagonal denominators divide $D$;
the remaining two divisibilities are exactly
\eqref{eq:fibrepolecertificate}. Thus no pole is inferred to
cancel merely from the rational module's rank.

Finally integrate the zero-block markers conditionally. Its
mean $tC_0B_0(u)$ and covariance $C_0$ are polynomial in $t,u$.
Every original marker becomes a polynomial in $t,u,Y_\lambda$.
The polynomial Gaussian means and covariances above make its
period polynomial as well. This proves safety.

For the stated bound, if the zero block has nilpotency index
$r_0$, its mean has $t$ degree at most $r_0\le s$ and covariance
at most $r_0-1$. Since $\deg D+r_0\le s$, the joint means
above have degree at most $s+1$ and covariances at most $s$.
If the zero block is absent, the same bounds hold directly.
Conditional Wick expansion of a marker of degree $q$ has
coefficient degree at most $sq$ in $t$ and remaining variable
degree at most $bq$. Its Gaussian expectation therefore has
degree at most $sq+(s+1)bq$, proving
\eqref{eq:fibrecutoff}.
\end{proof}

For a general matrix $A$, the rationality conclusion and the
identities remain valid. The unresolved finite-pole problem is
confined further: at each nonzero eigenvalue, purity forces its
largest Jordan block to equal the full algebraic multiplicity,
because that is the maximum possible pole order of $K$.
Thus there can be only one Jordan block there. We do not infer
the required marker-pole cancellation for these defective
nonzero blocks. The semisimple result and the nilpotent zero
block are complete; this last cancellation problem is separate.

\begin{thebibliography}{99}
\bibitem{Maisonobe}
P.~Maisonobe, \emph{Filtration relative, l'id\'eal de Bernstein et ses
pentes}, \emph{Rendiconti del Seminario Matematico della Universit\`a
di Padova} \textbf{150} (2023), 81--125;
\href{https://doi.org/10.4171/RSMUP/101}{doi:10.4171/RSMUP/101};
\href{https://arxiv.org/abs/1610.03354v1}{arXiv:1610.03354v1}.

\bibitem{SabbahStationary}
C.~Sabbah, \emph{An explicit stationary phase formula for the local
formal Fourier--Laplace transform}, in \emph{Singularities I},
Contemporary Mathematics \textbf{474}, American Mathematical Society,
2008, 309--330;
\href{https://arxiv.org/abs/0706.3570v2}{arXiv:0706.3570v2}.

\bibitem{Wall}
C.~T.~C.~Wall, \emph{Singular Points of Plane Curves}, Cambridge
University Press, 2004, Chapter~2, 15--38,
\href{https://doi.org/10.1017/CBO9780511617560.003}{doi:10.1017/CBO9780511617560.003}.

\end{thebibliography}
\end{document}
